% linear_congruences_solutions.tex
\documentclass[11pt]{article}
\usepackage{amsmath,amssymb}
\usepackage{geometry}
\geometry{margin=1in}
\title{Problems on Linear Congruences and Reduced Residue Systems}



% ------------------------------------------------------------------------------
% Required Packages for preamble/font.tex
% ------------------------------------------------------------------------------
\makeatletter
% Font encoding and families
\@ifpackageloaded{fontenc}{}{\usepackage[T1]{fontenc}}   % better font encoding
\@ifpackageloaded{lmodern}{}{\usepackage{lmodern}}       % Latin Modern — scalable replacement for CM

% Typography helpers
\@ifpackageloaded{lettrine}{}{\usepackage{lettrine}}      % dropped capitals
\@ifpackageloaded{lipsum}{}{\usepackage{lipsum}}          % dummy text (testing)
\makeatother

% ------------------------------------------------------------------------------
% Documentation and Usage Examples: lipsum and lettrine
% ------------------------------------------------------------------------------
% DESCRIPTION
%   The lipsum package provides dummy text for testing typographic layout.
%   The lettrine package produces dropped capitals at the beginning of a paragraph.
%
% FEATURES
%   - lipsum: generate one or more paragraphs for quick layout checks
%   - lettrine: control dropped capital lines, spacing, and scaling
%
% USAGE: lipsum
%   Basic:    \lipsum                  % generates a default paragraph
%   Paragraph: \lipsum[1]              % first paragraph
%   Range:     \lipsum[2-3]            % paragraphs 2 to 3
%   Words:     \lipsum[1][1-5]         % words 1–5 from paragraph 1
%
% USAGE: lettrine
%   Basic:     \lettrine{T}{his} starts a paragraph with a dropped “T”
%   Lines:     \lettrine[lines=3]{T}{his}    % drop across 3 lines
%   Spacing:   \lettrine[findent=2pt,nindent=0pt]{T}{his}
%   Scaling:   \lettrine[lines=3,loversize=0.15]{T}{his}
%   Slope:     \lettrine[hang=1]{T}{his}     % tighter hang of text under the drop cap
%
% EXAMPLES (copy into your document body to preview)
%   % Example 1: Simple lipsum
%   % \lipsum[1]
%
%   % Example 2: Variable Length lipsum
%   % \lipsum[2-3]
%
%   % Example 3: Drop cap with default settings
%   % \lettrine{L}{orem} ipsum dolor sit amet, consectetur adipiscing elit.\lipsum[2]
%
%   % Example 4: Three-line drop cap with adjusted size
%   % \lettrine[lines=3,loversize=0.2]{D}{olor} sit amet, consectetur adipiscing elit.\lipsum[3]
%
%   % Example 5: Fine-tuned indents for tighter text wrap
%   % \lettrine[lines=3,findent=1.5pt,nindent=0pt]{T}{ypography} testing block.\lipsum[4]
%
% NOTES
%   - Use \par after a lettrine paragraph if spacing looks off in complex layouts.
%   - Pair lettrine with a serif body font for best aesthetics.


% ------------------------------------------------------------------------------
% Required Packages for preamble/math.tex
% ------------------------------------------------------------------------------
\makeatletter
\@ifpackageloaded{amsmath}{}{\usepackage{amsmath}}      % (implicit) math foundations
\@ifpackageloaded{amssymb}{}{\usepackage{amssymb}}      % \mathfrak{...}, \square, \blacksquare
% \@ifpackageloaded{mathrsfs}{}{\usepackage{mathrsfs}}    % \mathscr{...}
% \@ifpackageloaded{amsfonts}{}{\usepackage{amsfonts}}    % math symbols + DeclarePairedDelimiter
% \@ifpackageloaded{mathtools}{}{\usepackage{mathtools}}  % enhanced math tools
% \@ifpackageloaded{dsfont}{}{\usepackage{dsfont}}        % \mathds{...}
% \@ifpackageloaded{xspace}{}{\usepackage{xspace}}        % text-usable macros
% \@ifpackageloaded{xparse}{}{\usepackage{xparse}}        % define flexible macros (\NewDocumentCommand)
\@ifpackageloaded{ifthen}{}{\usepackage{ifthen}}        % basic conditional logic (\ifthenelse)
\@ifpackageloaded{mdframed}{}{\usepackage{mdframed}}    % framed / boxed environments
\@ifpackageloaded{xcolor}{}{\usepackage{xcolor}         % color definitions and coloring text/math
  \definecolor{redcolor}{rgb}{1.0, 0.0, 0.0}            % {\color{redcolor}...}
  \definecolor{greencolor}{rgb}{0.3, 1.0, 0.3}          % {\color{greencolor}...}
  \definecolor{bluecolor}{rgb}{0.0, 0.0, 1.0}           % {\color{bluecolor}...}
  \definecolor{yellowcolor}{rgb}{1.0, 1.0, 0.0}         % {\color{yellowcolor}...}
  \definecolor{cyancolor}{rgb}{0.0, 1.0, 1.0}           % {\color{cyancolor}...}
  \definecolor{magentacolor}{rgb}{1.0, 0.0, 1.0}        % {\color{magentacolor}...}
  \definecolor{blackcolor}{rgb}{0, 0, 0}                % {\color{blackcolor}...}
  \definecolor{graycolor}{rgb}{0.2, 0.2, 0.2}           % {\color{graycolor}...}
  \definecolor{darkgreen}{RGB}{0,128,0}                 % {\color{darkgreen}...}
}
\makeatother












% ------------------------------------------------------------------------------
% \begin{define}[<color>]{<title>} … \end{define}
% ------------------------------------------------------------------------------
% DESCRIPTION:
%   This environment creates a framed definition or theorem with colored background
%   and border. It provides a visually distinct way to present important concepts.
%
% FEATURES:
%   - Colored frame and background (10% opacity)
%   - Section symbol (§) followed by title
%   - Configurable colors for customization
%   - Proper spacing and typography
%
% PARAMETERS:
%   [<color>]    Background and border color (optional, defaults to graycolor)
%   {<title>}    Title of the definition/theorem (required)
%
% DEPENDENCIES:
%   \usepackage{mdframed}
%   \usepackage{xcolor}
%
% USAGE EXAMPLES:
%   \begin{define}{Definition of Continuity}
%     A function f is continuous at a point a if...
%   \end{define}
%
%   \begin{define}[green]{Important Theorem}
%     This theorem is framed in green.
%   \end{define}
% ------------------------------------------------------------------------------
\newenvironment{define}[2][graycolor]
  {
  \vspace{2pt}
  \begin{mdframed}[backgroundcolor=#1!10, linecolor=#1]
  \color{#1}
  \textbf{\S~#2}\hspace{6pt}  % Start symbol for definition/theorem with a title
  \vspace{4pt}
  }
{
  \end{mdframed}
}

% ------------------------------------------------------------------------------
% \begin{claim}[<color>]{<title>} … \end{claim}
% ------------------------------------------------------------------------------
% DESCRIPTION:
%   This environment creates a framed claim or assertion with colored background
%   and border. It provides a visually distinct way to present important claims.
%
% FEATURES:
%   - Colored frame and background (10% opacity)
%   - Section symbol (§) followed by title
%   - Configurable colors for customization
%   - Proper spacing and typography
%
% PARAMETERS:
%   [<color>]    Background and border color (optional, defaults to bluecolor)
%   {<title>}    Title of the claim/assertion (required)
%
% DEPENDENCIES:
%   \usepackage{mdframed}
%   \usepackage{xcolor}
%
% USAGE EXAMPLES:
%   \begin{claim}{Key Claim}
%     This is an important claim that needs to be proven.
%   \end{claim}
%
%   \begin{claim}[red]{Critical Assertion}
%     This claim is framed in red.
%   \end{claim}
% ------------------------------------------------------------------------------
\newenvironment{claim}[2][bluecolor]
  {
  \vspace{2pt}
  \begin{mdframed}[backgroundcolor=#1!10, linecolor=#1]
  \color{#1}
  \textbf{\S~#2}\hspace{6pt}  % Start symbol for claim/assertion with a title
  \vspace{4pt}
  }
{
  \end{mdframed}
}



% ------------------------------------------------------------------------------
% \begin{state}[<number>]{<title>} ... \end{state}
% ------------------------------------------------------------------------------
% DESCRIPTION:
%   This environment creates a numbered subsection title with automatic counter
%   management. It formats subsection titles in the style "§<number> (<title>)".
%
% FEATURES:
%   - Automatic counter incrementing
%   - Optional counter reset to specific value
%   - Consistent formatting with section symbol (§)
%   - Unnumbered subsection with custom title format
%
% PARAMETERS:
%   [<number>]    Counter reset value (optional, resets counter to this value)
%   {<title>}     Title text for the subsection (required)
%
% DEPENDENCIES:
%   \usepackage{xparse}
%
% USAGE EXAMPLES:
%   \begin{state}{First Section}
%     Content of the first section...
%   \end{state}
%
%   \begin{state}[5]{Reset Section}
%     Content with counter reset to 5...
%   \end{state}
% ------------------------------------------------------------------------------
\newcounter{statecounter}
\NewDocumentEnvironment{state}{ o m }{%
  \IfNoValueTF{#1}{%
    \stepcounter{statecounter}%
  }{%
    \setcounter{statecounter}{\numexpr#1-1\relax}%
    \stepcounter{statecounter}%
  }%
  \subsection*{§\arabic{statecounter} (#2)}%
}{}

% ------------------------------------------------------------------------------
% Alias environment "prop" as a wrapper for "state".
% The following two commands are identical:
% \begin{state} ... \end{state}
% \begin{prop} ... \end{prop}
% ------------------------------------------------------------------------------
\NewDocumentEnvironment{prop}{ o m }{%
  % Pass the arguments to the "state" environment.
  \IfNoValueTF{#1}{%
    \begin{state}{#2}%
  }{%
    \begin{state}[#1]{#2}%
  }%
}{%
  \end{state}%
}

% ------------------------------------------------------------------------------
% \begin{subanswer}[<Color>][<isWhite>] … \end{subanswer}
% ------------------------------------------------------------------------------
% DESCRIPTION:
%   This custom environment displays an "Answer:" label with configurable text color.
%   It is intended for formatting answers clearly in assignments or solutions.
%
% FEATURES:
%   - Displays "Answer:" in bold, followed by answer content.
%   - Automatically ends with a square symbol (□) aligned right.
%   - First optional argument selects text color (default = red).
%   - Second optional argument overrides to white text when set to true.
%
% PARAMETERS (both optional):
%   [<Color>]    Text color name (e.g. blue, green). Defaults to red.
%   [<isWhite>]  true  -> force white text (for dark backgrounds)
%                 false -> use the chosen <Color>. Defaults to false.
%
% DEPENDENCIES:
%   \usepackage{ifthen}
%   \usepackage{xcolor}
%   \usepackage{xparse}
%
% USAGE EXAMPLES:
%   % defaults: red text, isWhite = false
%   \begin{subanswer}
%     This answer is red.
%   \end{subanswer}
%
%   % blue text, isWhite = false
%   \begin{subanswer}[blue]
%     This answer is blue.
%   \end{subanswer}
%
%   % green text, forced white override
%   \begin{subanswer}[green][true]
%     This answer is white (on green background).
%   \end{subanswer}
% ------------------------------------------------------------------------------
\NewDocumentEnvironment{subanswer}{ O{red} O{false} }
{%
  \par\medskip                     % start a paragraph with a bit of space
  % choose white if second optional arg is "true", otherwise #1
  \ifthenelse{\equal{#2}{true}}%
    {\color{white}}%
    {\color{#1}}%
  \noindent\textbf{Answer}:        % no \\ or \newline here
}
{%
  \par\hfill$\square$\par
}

% ------------------------------------------------------------------------------
% \begin{subprove}[<Color>][<isWhite>] … \end{subprove}
% ------------------------------------------------------------------------------
% DESCRIPTION:
%   This custom environment displays a "Proof:" label with configurable text color.
%   It is intended for formatting mathematical proofs clearly in assignments or solutions.
%
% FEATURES:
%   - Displays "Proof:" in bold, followed by proof content.
%   - Automatically ends with a filled square symbol (■) aligned right.
%   - First optional argument selects text color (default = red).
%   - Second optional argument overrides to white text when set to true.
%
% PARAMETERS (both optional):
%   [<Color>]    Text color name (e.g. blue, green). Defaults to red.
%   [<isWhite>]  true  -> force white text (for dark backgrounds)
%                 false -> use the chosen <Color>. Defaults to false.
%
% DEPENDENCIES:
%   \usepackage{ifthen}
%   \usepackage{xcolor}
%   \usepackage{xparse}
%
% USAGE EXAMPLES:
%   % defaults: red text, isWhite = false
%   \begin{subprove}
%     This proof is red.
%   \end{subprove}
%
%   % blue text, isWhite = false
%   \begin{subprove}[blue]
%     This proof is blue.
%   \end{subprove}
%
%   % green text, forced white override
%   \begin{subprove}[green][true]
%     This proof is white (on green background).
%   \end{subprove}
% ------------------------------------------------------------------------------
\NewDocumentEnvironment{subprove}{ O{red} O{false} }
{%
  \par\medskip                     % start a paragraph with a bit of space
  % choose white if second optional arg is "true", otherwise #1
  \ifthenelse{\equal{#2}{true}}%
    {\color{white}}%
    {\color{#1}}%
  \noindent\textbf{Proof}:        % no \\ or \newline here
}
{%
  \par\hfill$\blacksquare$\par
}

% ------------------------------------------------------------------------------
% \begin{answer}[<color>] … \end{answer}
% ------------------------------------------------------------------------------
% DESCRIPTION:
%   This environment creates an answer section with optional text color.
%   It automatically adds a new page after the answer for clean separation.
%
% FEATURES:
%   - Displays "Answer:" in bold with specified color
%   - Automatically ends with a square symbol (□) aligned right
%   - Adds a new page after the answer
%   - Optional color parameter for customization
%
% PARAMETERS:
%   [<color>]    Text color name (optional, defaults to redcolor)
%
% DEPENDENCIES:
%   None (uses basic LaTeX commands)
%
% USAGE EXAMPLES:
%   \begin{answer}
%     This is a default red answer.
%   \end{answer}
%
%   \begin{answer}[blue]
%     This is a blue answer.
%   \end{answer}
% ------------------------------------------------------------------------------
\newenvironment{answer}[1][redcolor]
  {\ \newline \color{#1}\noindent\textbf{Answer}: }
  {\par\hfill$\square$\newline\newpage}

% ------------------------------------------------------------------------------
% \begin{prove}[<color>] … \end{prove}
% ------------------------------------------------------------------------------
% DESCRIPTION:
%   This environment creates a proof section with optional text color.
%   It automatically adds a new page after the proof for clean separation.
%
% FEATURES:
%   - Displays "Proof:" in bold with specified color
%   - Automatically ends with a filled square symbol (■) aligned right
%   - Adds a new page after the proof
%   - Optional color parameter for customization
%
% PARAMETERS:
%   [<color>]    Text color name (optional, defaults to redcolor)
%
% DEPENDENCIES:
%   None (uses basic LaTeX commands)
%
% USAGE EXAMPLES:
%   \begin{prove}
%     This is a default red proof.
%   \end{prove}
%
%   \begin{prove}[blue]
%     This is a blue proof.
%   \end{prove}
% ------------------------------------------------------------------------------
\newenvironment{prove}[1][redcolor]
  {\ \newline \color{#1}\noindent\textbf{Proof}: }
  {\par\hfill$\blacksquare$\newline\newpage}


\input{../preamble/operator.tex}
\input{../preamble/symbol.tex}
%%%%%%%%%%%%%%%%%%%%%%%%%%%%%%%%%%%%%%%%%%%%%%%%%%%%
% Grouping Shortcuts for VS Code (keybindings.json)
%%%%%%%%%%%%%%%%%%%%%%%%%%%%%%%%%%%%%%%%%%%%%%%%%%%%
% {
%   "key": "ctrl+shift+9",
%   "command": "editor.action.insertSnippet",
%   "when": "editorTextFocus && editorLangId == 'latex'",
%   "args": {
%     "snippet": "\\group{${TM_SELECTED_TEXT:${1}}}"
%   }
% },
% {
%   "key": "ctrl+9",
%   "command": "editor.action.insertSnippet",
%   "when": "editorTextFocus && editorLangId == 'latex'",
%   "args": {
%     "snippet": "\\arr{${TM_SELECTED_TEXT:${1}}}"
%   }
% },
% {
%   "key": "ctrl+shift+[",
%   "command": "editor.action.insertSnippet",
%   "when": "editorTextFocus && editorLangId == 'latex'",
%   "args": {
%     "snippet": "\\set{${TM_SELECTED_TEXT:${1}}}"
%   }
% },
% {
%   "key": "ctrl+shift+\\",
%   "command": "editor.action.insertSnippet",
%   "when": "editorTextFocus && editorLangId == 'latex'",
%   "args": {
%     "snippet": "\\abs{${TM_SELECTED_TEXT:${1}}}"
%   }
% },
% {
%   "key": "ctrl+shift+,",
%   "command": "editor.action.insertSnippet",
%   "when": "editorTextFocus && editorLangId == 'latex'",
%   "args": {
%     "snippet": "\\coord{${TM_SELECTED_TEXT:${1}}}"
%   }
% },

% ------------------------------------------------------------------------------
% Required Packages for preamble/group.tex
% ------------------------------------------------------------------------------
\makeatletter
\@ifpackageloaded{amsmath}{}{\usepackage{amsmath}}      % (implicit) math foundations
\@ifpackageloaded{xspace}{}{\usepackage{xspace}}        % text-usable macros
% \@ifpackageloaded{dsfont}{}{\usepackage{dsfont}}        % \mathds{...}
% \@ifpackageloaded{mathrsfs}{}{\usepackage{mathrsfs}}    % \mathscr{...}
% \@ifpackageloaded{amssymb}{}{\usepackage{amssymb}}      % \mathfrak{...}
% \@ifpackageloaded{amsfonts}{}{\usepackage{amsfonts}}    % math symbols + DeclarePairedDelimiter
% \@ifpackageloaded{mathtools}{}{\usepackage{mathtools}}  % enhanced math tools
\makeatother



% ------------------------------------------------------------------------------
% Mathematical Grouping Commands
% ------------------------------------------------------------------------------
% DESCRIPTION:
%   Commands for mathematical grouping and delimiters with automatic sizing.
%   These commands are designed for quick insertion via VS Code keybindings.
%
% FEATURES:
%   - Automatic sizing with \left and \right
%   - Safe to use in text mode with \ensuremath and \xspace
%   - VS Code keybinding integration
%   - Consistent formatting across documents
%
% DEPENDENCIES:
%   \usepackage{xspace}
%
% USAGE EXAMPLES:
%   The inner product \coord{x,y} defines a metric.
%   The group \group{G} acts on the set \set{X}.
%   The absolute value \abs{x} is always non-negative.
%   The interval \arr{a,b} contains all real numbers between a and b.
%   The inner product $\coord{x,y}$ defines a metric.
%   The group $\group{G}$ acts on the set $\set{X}$.
%   The absolute value $\abs{x}$ is always non-negative.
%   The interval $\arr{a,b}$ contains all real numbers between a and b.
% ------------------------------------------------------------------------------
\newcommand{\group}[1]{\ensuremath{\left(#1\right)}\xspace}              % Parentheses: (x)
\newcommand{\arr}[1]{\ensuremath{\left[#1\right]}\xspace}                % Square brackets: [x]
\newcommand{\set}[1]{\ensuremath{\left\{#1\right\}}\xspace}              % Curly braces: {x}
\newcommand{\abs}[1]{\ensuremath{\left|#1\right|}\xspace}                % Absolute value: |x|
\newcommand{\coord}[1]{\ensuremath{\left\langle#1\right\rangle}\xspace}  % Angle brackets: ⟨x⟩





\begin{document}


\begin{define}{Divisibility}[thm:divisibility][red]
  For integers $a,b$ and positive integer $n$, $n\mid a$ if and only if there exists an integer $k$ such that
  \[
    a = kn.
  \]
\end{define}

\begin{define}{Coprime (Relatively Prime)}[def:coprime][red]
  Two integers $a$ and $b$ are \textbf{coprime} if $\gcd(a,b)=1$.
\end{define}

\begin{define}{Greatest Common Divisor (gcd)}[def:gcd][red]
  For integers $a,b$, not both zero, the greatest common divisor of $a$ and $b$, denoted $\gcd(a,b)$, is the largest positive integer that divides both $a$ and $b$.
\end{define}

\begin{enumerate}
\item Integers $a,b$ with $\gcd(a,b)=3$ and $a+b=65$.
\begin{subanswer}
  If $\gcd(a,b)=3$, then $3\mid a$ and $3\mid b$, so $3\mid(a+b)$. But $3\nmid65$. Contradiction. \textbf{No such integers.}
\end{subanswer}
\end{enumerate}


\section*{Linear Congruences}

\begin{define}{Congruence}[thm:congruence-properties][red]
  For integers $a,b$ and positive integer $n$, $n\mid(a-b)$ if
  \[
    a\equiv b\pmod n,
  \]
  where $a$ and $b$ have the same \textbf{residue} (\textbf{remainder}) when divided by $n$. 
  It follows that:
  \begin{itemize}
    \item If $a\equiv b\pmod n$ and $c\equiv d\pmod n$, then $a+c\equiv b+d\pmod n$.
    \item If $a\equiv b\pmod n$ and $c\equiv d\pmod n$, then $ac\equiv bd\pmod n$.
    \item If $a\equiv b\pmod n$ and $d\mid n$, then $a\equiv b\pmod d$.
  \end{itemize}
  \textbf{Solvability:} Let $d=\gcd(a,n)$. The congruence $ax\equiv b\pmod n$ has $d$ distinct solutions modulo $n$ iff 
  $$d\mid b$$
  with the least positive solution $x$ where $(a/d)x\equiv (b/d)\pmod{n/d}$.
\end{define}


\begin{enumerate}
  \item $84x\equiv5\pmod{49}$.
  \begin{subanswer}
    $\gcd(84,49)=7$. Since $7\nmid5$, by theorem 1 there is \textbf{no solution}.
  \end{subanswer}
  \item $12x\equiv16\pmod{32}$.
  \begin{subanswer}
    $\gcd(12,32)=4$ and $4\mid16$. Divide by $4$: $3x\equiv4\pmod8$. Since $\gcd(3,8)=1$ and $3^{-1}\equiv3\pmod8$, we get $x\equiv3\cdot4=12\equiv4\pmod8$. The least positive solution is \boxed{4}. (All solutions mod $32$: $4,12,20,28$.)
  \end{subanswer}
\end{enumerate}






\newpage

\section*{Multiplicative Inverses in Congruences}

\begin{define}{Multiplicative Inverse Modulo $n$}[thm:inverse][red]
  An integer $a$ has a multiplicative inverse modulo $n$ if and only if $\gcd(a,n)=1$.
\end{define}

\begin{enumerate}
  \item Find the inverse of $7$ modulo $24$ (i.e., solve $7x\equiv1\pmod{24}$).
  \begin{subanswer}
    Since $\gcd(7,24)=1$, an inverse exists. Use the extended Euclidean algorithm:
    \begin{align*}
      \underbar{24} &= 3 \cdot \underbar{7} + 3 \\
      7 &= 2 \cdot 3 + \underbar{1} \\
      3 &= 3 \cdot 1 + 0
    \end{align*}
    Back-substitute, $(7) \cdot \underbar{7} + (-2) \cdot \underbar{24}=\underbar{1}.$
    That is, $$\underbrace{7}_{x} \cdot \underbar{7} = 2 \cdot \underbar{24} + \underbar{1}.$$
    Thus, $7^{-1} \equiv 7 \pmod{24}$. Least positive solution: $\boxed{7}$.
  \end{subanswer}

  \item Find the inverse of $831$ modulo $402$ (i.e., solve $831x\equiv1\pmod{402}$) or show that no such inverse exists.
  \begin{subanswer}
    By the extended Euclidean algorithm:
    \begin{align*}
      \underbar{831} &= 2 \cdot \underbar{402} + 27 \\
      402 &= 14 \cdot 27 + 24 \\
      27 &= 1 \cdot 24 + \underbar{3} \\
      24 &= 8 \cdot 3 + 0
    \end{align*}
    Back-substitute, $(15) \cdot \underbar{831} + (- 31) \cdot \underbar{402}=\underbar{3}.$ Since $\gcd(831,402)=3\neq1$, \textbf{no inverse exists}.
    But still, $$\underbrace{15}_{x} \cdot \underbar{831} = (31) \cdot \underbar{402} + \underbar{3}.$$
    That is, $831x\equiv3\pmod{402}$ has solutions. 
    Dividing by $3$, $277x\equiv1\pmod{134}$. Since $\gcd(277,134)=1$, an inverse exists. By the extended Euclidean algorithm,
    \begin{align*}
      \underbar{277} &= 2 \cdot \underbar{134} + 9 \\
      134 &= 14 \cdot 9 + 8 \\
      9 &= 1 \cdot 8 + \underbar{1} \\
      8 &= 8 \cdot 1 + 0
    \end{align*}
    Back-substitute, $(15) \cdot \underbar{277} + (- 31) \cdot \underbar{134}=\underbar{1}.$
    That is, $$\underbrace{15}_{x} \cdot \underbar{277} = 31 \cdot \underbar{134} + \underbar{1}.$$
    Thus, $277^{-1} \equiv 15 \pmod{134}$. Least positive solution: $\boxed{15}$. (All solutions mod $402$: $15,149,283$.)
    
  \end{subanswer}
\end{enumerate}







\newpage

\section*{Systems of Linear Congruences}


\begin{define}{Chinese Remainder Theorem}[thm:crt][redcolor]
  Let $m_1, m_2, \dots, m_n$ be pairwise coprime positive integers, and let $b_1, b_2, \dots, b_n \in \mathbb{Z}$. Then the system of congruences
  \[
    \begin{cases}
      x \equiv b_1 \pmod{m_1} \\
      x \equiv b_2 \pmod{m_2} \\
      \;\;\vdots \\
      x \equiv b_n \pmod{m_n}
    \end{cases}
  \]
  has a solution $x \in \mathbb{Z}$, and all solutions are congruent modulo $M := m_1 m_2 \cdots m_n$.
\end{define}

\begin{enumerate}

  \item System $x\equiv3\pmod7$, $x\equiv4\pmod6$.
  \begin{subanswer}
    $\gcd(7,6)=1$. Let $x=3+7k$. Then $3+7k\equiv4\pmod6\Rightarrow7k\equiv1\pmod6$. Since $7\equiv1\pmod6$, $k\equiv1\pmod6$. Take $k=1$ gives $x=10$. Hence $x\equiv10\pmod{42}$. Least positive: \boxed{10}.
  \end{subanswer}
  
  \item 
  \begin{align*}
  x & \equiv 3 \pmod{4} \\
  x & \equiv 2 \pmod{5}
  \end{align*}
  \begin{subanswer}
    Since $\gcd(4,5)=1$, there is a unique solution modulo $20$. Let $x=4k+3$. 
    Then $4k+3\equiv 2\pmod{5}$ gives $4k\equiv 4\pmod{5}$, so $k\equiv 1\pmod{5}$ 
    and $x=20m+7$.\vspace*{10px}\\
    Therefore, by the \nameref{thm:crt}, the least non-negative solution is $x=7$.
  \end{subanswer}

  \item 
  \begin{align*}
  3x & \equiv 2 \pmod{4} \\
  4x & \equiv 1 \pmod{5} \\
  6x & \equiv 3 \pmod{9}
  \end{align*}
  \begin{subanswer}
    Reduce each congruence. Since $3^{-1}\equiv 3\pmod{4}$, $x\equiv 2\cdot 3\equiv 2\pmod{4}$. 
    Since $4^{-1}\equiv 4\pmod{5}$, $x\equiv 1\cdot 4\equiv 4\pmod{5}$. That is,
    \begin{align*}
    x & \equiv 2 \pmod{4} \\
    x & \equiv 4 \pmod{5} \\
    x & \equiv 2 \pmod{3}.
    \end{align*}
    Let $x=4k+2$. Then $4k\equiv 0\pmod{3}$ so $k\equiv 0\pmod{3}$, giving $x\equiv 2\pmod{12}$. 
    With $x=12k+2$ and $x\equiv 4\pmod{5}$, we obtain $12k\equiv 2\pmod{5}$, 
    so $k\equiv 1\pmod{5}$ and $x=60n+14$.\vspace*{10px}\\
    Therefore, by the \nameref{thm:crt}, the least non-negative solution is $x=14$.
  \end{subanswer}
\end{enumerate}




\newpage




\section*{Reduced Residue Systems}
\begin{define}{Reduced Residue System Modulo $n$}[][red]
  The symbol $(\mathbb{Z}/n\mathbb{Z})^\times$ denotes the 
  \textbf{reduced residue system modulo $n$}, also called the 
  \textbf{group of units modulo $n$}. It consists of all 
  \textbf{congruence classes} $[a]$ in $\mathbb{Z}/n\mathbb{Z}$ 
  such that $\gcd(a, n) = 1$:
  \[
    (\mathbb{Z}/n\mathbb{Z})^\times = \{ [a] \in \mathbb{Z}/n\mathbb{Z} \mid \gcd(a, n) = 1 \}.
  \]
  where 
  $$[a]=\{ x \in \mathbb{Z} \mid x\equiv a \pmod{n} \}$$ 
  is the equivalence class of integers congruent to $a$ modulo $n$. It follows that
  \begin{itemize}
    \item \textbf{Closure:} If $[a], [b] \in (\mathbb{Z}/n\mathbb{Z})^\times$, then $[a][b] \in (\mathbb{Z}/n\mathbb{Z})^\times$.
    \item \textbf{Identity:} $[1]$ is the multiplicative identity.
    \item \textbf{Inverse:} Each $[a]$ has an inverse $[b]$ such that $[a][b] = [1]$.
    \item \textbf{Associativity:} Multiplication modulo $n$ is associative.
  \end{itemize}
\end{define}
\begin{enumerate}
  \item List the \textbf{least non-negative reduced residue system} modulo $10$.
  \begin{subanswer}
    A reduced residue system modulo $10$ is a set of $\varphi(10)=4$ integers that are coprime to $10$. That is,
    $\gcd(1, 10) = 1$, $\gcd(3, 10) = 1$, $\gcd(7, 10) = 1$, and $\gcd(9, 10) = 1$. 
    Thus, the units (least non-negative) reduced residue system is $\boxed{\{1, 3, 7, 9\}}$.
  \end{subanswer}

  \item Is $\{-15,66,-75,11\}$ a reduced residue system modulo $8$?
  \begin{subanswer}
    Reduce mod $8$: $-15\equiv1$, $66\equiv2$, $-75\equiv5$, $11\equiv3$. Residues $\{1,2,5,3\}$. But $2$ is not coprime to $8$ (gcd $=2$). Thus the set is \textbf{not} a reduced residue system modulo $8$.
  \end{subanswer}

  \item Reduced residue system modulo $7$ consisting only of odd integers.
  \begin{subanswer}
    One choice: \boxed{\{1,3,5,9,11,13\}}. Modulo $7$ these reduce to $\{1,3,5,2,4,6\}$, i.e. all units.
  \end{subanswer}
\end{enumerate}



\section*{Factorials in Congruences}


\begin{define}{Wilson's Theorem}[thm:wilson][red]
  For prime $p$, $(p-1)!\equiv-1\pmod p$.
\end{define}


\begin{enumerate}
\item For odd prime $p$, prove $2(p-3)!\equiv-1\pmod p$.
\begin{subprove}
  By Wilson's theorem, $(p-1)! \equiv -1 \pmod p$. But $(p-1)! = (p-1)(p-2)(p-3)! \equiv (-1)(-2)(p-3)! \equiv 2(p-3)! \pmod p$. Thus $2(p-3)! \equiv -1 \pmod p$.
\end{subprove}
\end{enumerate}


\newpage





\section*{Multiplicative Arithmetic Functions}


\begin{define}{Multiplicativity}[def:multiplicativity][red]
  A function $f:\mathbb{Z}_{>0}\to\mathbb{C}$ is \textbf{multiplicative} if 
  \[
    f(mn) = f(m)f(n) \text{ for any coprime positive integers } m,n.
  \]
  \textbf{completely multiplicative} if
  \[
    f(mn) = f(m)f(n) \text{ for all positive integers } m,n.
  \]
\end{define}


\begin{define}{Chinese Remainder Theorem (pair):}[thm:CRT-pair][red]
  If $\gcd(m,n)=1$ (or coprime), then
  \[(\mathbb{Z}/mn\mathbb{Z})^\times\cong(\mathbb{Z}/m\mathbb{Z})^\times\times(\mathbb{Z}/n\mathbb{Z})^\times.
  \]
  Taking cardinalities gives multiplicativity of $\varphi$ on coprime arguments.
\end{define}


\begin{define}{Euler's Totient Function $\varphi(n)$}[def:phi][red]
  The function $\varphi:\mathbb{Z}_{>0}\to\mathbb{Z}_{>0}$ is defined by
  \[
    \varphi(n)=|(\mathbb{Z}/n\mathbb{Z})^\times| = \abs{\sBUILD{1 \leq k \leq n}{\group{k,n} = 1}},
  \]
  i.e. $\varphi(n)$ is the number of integers in $\{1,2,\dots,n\}$ 
  that are coprime to $n$.\\
  \textbf{Note:} If $n$ is prime, then $\phi(n) = n-1$.
\end{define}


\begin{enumerate}
  \item Prove $\varphi(p^k)=p^k-p^{k-1}=p^{k-1}(p-1)$ for prime $p$ and integer $k\ge1$.
  \begin{subprove}
    The integers in $\{1,2,\dots,p^k\}$ that are not coprime to $p^k$ are exactly the 
    multiples of $p$: $1p,2p,3p,\dots,p^{k-1}p$. There are $p^{k-1}$ such multiples. Thus,
    \[
      \varphi(p^k) = p^k - p^{k-1} = p^{k-1}(p-1).
    \]
  \end{subprove}
  \item Prove that $\varphi$ is multiplicative but not completely multiplicative.
  \begin{subprove}
    If $\gcd(m,n)=1$ then by the Chinese Remainder Theorem,
    \[
      |(\mathbb{Z}/mn\mathbb{Z})^\times| = |(\mathbb{Z}/m\mathbb{Z})^\times| \cdot |(\mathbb{Z}/n\mathbb{Z})^\times|,
    \]
    so $\varphi(mn) = \varphi(m)\varphi(n)$. Thus, $\varphi$ is \textbf{multiplicative}.
    
    However, for prime $p$, $\varphi(p^k)=p^k-p^{k-1}=p^{k-1}(p-1)$. Let $k=2$. Then $\varphi(p^2)=p(p-1)$ but 
    $\varphi(p)\varphi(p)=(p-1)(p-1)=(p-1)^2\neq p(p-1)$ for $p>2$. 
    Thus, $\varphi$ is \textbf{not completely multiplicative}.
  \end{subprove}
\end{enumerate}




















\newpage



























\newpage
\section*{Exponentiation in Congruences}




\begin{define}{Euler's theorem}[thm:euler][red]
  If $\gcd(a,m)=1$ then $$a^{\varphi(m)}\equiv1\pmod m.$$
  \begin{subprove}
    Let $R = \{ r_1, r_2, \dots, r_{\phi(n)} \}$ be the reduced residue system modulo $n$.
    Since $\gcd(a,n) = 1$, the products $a r_1, a r_2, \dots, a r_{\phi(n)}$ form another reduced residue system modulo $n$.
    Thus,
    \[
      a r_1 \cdot a r_2 \cdots a r_{\phi(n)} \equiv r_1 r_2 \cdots r_{\phi(n)} \pmod{n}.
    \]
    That is,
    \[
      a^{\phi(n)} (r_1 r_2 \cdots r_{\phi(n)}) \equiv r_1 r_2 \cdots r_{\phi(n)} \pmod{n}.
    \]
    Since $r_1 r_2 \cdots r_{\phi(n)}$ are coprimes of $n$, 
    \[
      a^{\phi(n)} \equiv 1 \pmod{n}.
    \]
  \end{subprove}
  \noindent\textbf{Corollary: Fermat's little theorem:} If $\gcd(a,m)=1$ and \underbar{$m$ is prime} then $$a^{m-1}\equiv1\pmod m.$$
\end{define}



\begin{enumerate}
  \item If $\gcd(a,m)=1$, show $a^{\varphi(m)-1}$ is the inverse of $a$ modulo $m$
  \begin{subprove}
    By Euler's theorem $a^{\varphi(m)}\equiv1\pmod m$. Since $a^{\varphi(m)}=a\cdot a^{\varphi(m)-1}$,
    \[a\cdot a^{\varphi(m)-1}\equiv a^{\varphi(m)}\equiv1\pmod m,
    \]
    so $a^{\varphi(m)-1}$ is an inverse of $a$ modulo $m$. (Cited: Euler's theorem.)
  \end{subprove}
  \item Find the least positive solution $x$ to the congruence $x \equiv 20^{110}(\bmod 17)$. 
    \begin{subanswer}
      $x\equiv20^{110}\pmod{17}$. Reduce base: $20\equiv3\pmod{17}$. Since $17$ is prime, by Fermat, $3^{16}\equiv1\pmod{17}$, and $110\equiv14\pmod{16}$. Thus
      \[3^{110}\equiv3^{14}=(3^8)(3^4)(3^2)\equiv16\cdot13\cdot9\equiv4\cdot9\equiv36\equiv2\pmod{17}.
      \]
      So least positive solution: $\boxed{2}$.
    \end{subanswer}
  \item For $x \equiv 38^{110}(\bmod 21)$.
    \begin{subanswer}
      $x\equiv38^{110}\pmod{21}$. Reduce base: $38\equiv17\pmod{21}$. 
      By Euler theorem and Multiplicativity of $\varphi$, $\varphi(21)=\varphi(3)\varphi(7)=2\cdot6=12$, 
      $$17^{\varphi(21)}\equiv17^{12}\equiv1\pmod{21}.$$
      Thus,
      \[38^{110}\equiv17^{110}\equiv17^2\equiv289\equiv16\pmod{21}.
      \]
      So least positive solution: $\boxed{16}$.
    \end{subanswer}
\end{enumerate}






\newpage





\begin{enumerate}
  \item Prove Euler's theorem.
  \begin{subprove}
    Let $m>1$ and let $\{r_1,\dots,r_{\varphi(m)}\}$ be a reduced residue system modulo $m$. Since $\gcd(a,m)=1$, the set $\{ar_1,\dots,ar_{\varphi(m)}\}$ is also a reduced residue system modulo $m$ (by problem 1 in next section). Thus,
    \[
      r_1 r_2 \cdots r_{\varphi(m)} \equiv (ar_1)(ar_2) \cdots (ar_{\varphi(m)}) \equiv a^{\varphi(m)}(r_1 r_2 \cdots r_{\varphi(m)}) \pmod m.
    \]
    Since $\gcd(r_1 r_2 \cdots r_{\varphi(m)}, m) = 1$, we can cancel it from both sides to get
    \[
      a^{\varphi(m)} \equiv 1 \pmod m.
    \]
  \end{subprove}

  \item Let $m>1$ and let $\{r_1,\dots,r_{\varphi(m)}\}$ be a reduced residue 
  system modulo $m$. Characterize integers $j$ for which 
  $\{jr_1,\dots,jr_{\varphi(m)}\}$ is also a reduced residue system modulo $m$.
  \begin{subprove}
    If $\gcd(j,m)=1$ then multiplication by $j$ is a bijection on $(\mathbb Z/m\mathbb Z)^\times$, so the set $\{jr_i\}$ consists of distinct residues all coprime to $m$, hence is a reduced residue system. Conversely, if $\gcd(j,m)>1$ then any $r_i$ with $\gcd(r_i,m)=1$ yields $\gcd(jr_i,m)>1$, so the image contains nonunits and cannot be a reduced residue system.
  \end{subprove}





  \item If $a,b\in\mathbb Z$, $a\ge b>0$, and $a=bq+r$ (with $0\le r<b$), show $(a,b)=(b,r)$.
  \begin{subprove}
    Let $d=(a,b)$. Then $d\mid a$ and $d\mid b$, so $d\mid(a-bq)=r$. Thus, $d\mid b$ and $d\mid r$, so $d\le(b,r)$. Conversely, let $e=(b,r)$. Then $e\mid b$ and $e\mid r$, so $e\mid(bq+r)=a$. Thus, $e\mid a$ and $e\mid b$, so $e\le(a,b)$. Therefore, $(a,b)=(b,r)$.
  \end{subprove}

  \item If $a,b\in\mathbb Z$ and $(a,b)=[a,b]$, what can you say about $a,b$?
  \begin{subprove}
    If $(a,b)=[a,b]$ then $a$ and $b$ are coprime, i.e., $\gcd(a,b)=1$.
  \end{subprove}

  \item Let $a,b\in\mathbb Z$ and $p$ prime. If $p\mid ab$ then $p\mid a$ or $p\mid b$.
  \begin{subprove}
    If $p\nmid a$ then $\gcd(a,p)=1$, so there exist integers $x,y$ such that $ax+py=1$. Multiplying by $b$ gives $abx+pby=b$. Since $p\mid ab$ and $p\mid pby$, we have $p\mid b$.

  \end{subprove}

  \item Let $p$ be prime and $p\nmid a$. Show there is a unique solution to $ax\equiv1\pmod p$.
  \begin{subprove}
    Since $\gcd(a,p)=1$, $a$ is a unit in $\mathbb Z/p\mathbb Z$ so an inverse $x$ exists. If $x_1,x_2$ are both inverses then $a(x_1-x_2)\equiv0\pmod p$. As $p\nmid a$, we get $p\mid(x_1-x_2)$, so $x_1\equiv x_2\pmod p$. Thus the inverse is unique modulo $p$.
  \end{subprove}

  \item Let $a,b,c,m\in\mathbb Z$ with $m,c\ge1$. If $a\equiv b\pmod m$ then $ac\equiv bc\pmod{mc}$.
  \begin{subprove}
    $a\equiv b\pmod m$ implies $m\mid(a-b)$, so $mc\mid c(a-b)=ac-bc$. Therefore $ac\equiv bc\pmod{mc}$.
  \end{subprove}

  \item Let $a,b,c,d\in\mathbb Z$ with $b,d\neq0$ and $(a,b)=(c,d)=1$. If $\dfrac{a}{b}+\dfrac{c}{d}$ is an integer, show $b=\pm d$.
  \begin{subprove}
    Suppose $\dfrac{a}{b}+\dfrac{c}{d}=k\in\mathbb Z$. Clearing denominators gives
    \[ad+bc=kbd\quad\Rightarrow\quad ad=b(kd-c).\]
    Thus $b\mid ad$. Since $\gcd(a,b)=1$, we conclude $b\mid d$. Symmetrically $d\mid b$. Hence $|b|=|d|$, so $b=\pm d$.
  \end{subprove}

\end{enumerate}





\end{document}
